\chapter{Couverture profonde et apprentissage automatique dans la tarification}\label{ch:ml:deep-hedging-and-machine-learning-in-pricing}

Une option d'achat à un mois à la monnaie, sur une action dont la volatilité suit le modèle de Heston, vaut 2.25 pour un spot de 100, et un point
de volatilité implicite en vaut 0.115 de sa valeur. Couverte chaque jour jusqu'à son delta de Black--Scholes, avec un coût de transaction de 5 points de base,
la couverture coûte en moyenne 0.123 en frais, soit plus d'un point de volatilité, pour un trader dont l'edge est souvent plus faible. Un
réseau qui apprend quand ne pas trader paie 0.099 et porte moins de risque. Ce chapitre utilise l'apprentissage automatique là où
les desks de produits dérivés l'utilisent~: pour couvrir en présence de frictions que le manuel ignore (deep hedging), pour remplacer un pricer lent par une
approximation rapide (\omterm{def:ml:deep-hedging-and-machine-learning-in-pricing:surrogate}{modèles de substitution}, ou surrogates, et \omterm{def:ml:deep-hedging-and-machine-learning-in-pricing:surrogate}{apprentissage automatique différentiel}, qui les rend précis avec peu d'échantillons), et pour calibrer
un modèle en une seule passe. Chacun est mesuré par rapport à la réponse exacte que fournissent les moteurs du livre~5, et chacun a un domaine où il
échoue.

\section{Couverture sous frictions comme problème d'apprentissage}

Le livre~5 (chapitre~4) a mesuré l'erreur de couverture discrète du delta de Black--Scholes et (chapitre~26) la bande de couverture que
les coûts de transaction rendent nécessaire. Avec des coûts, la meilleure couverture dépend de la position détenue (revenir au delta ne vaut pas la peine
si le delta a peu bougé), du risque que le trader est prêt à porter, et de tout ce que le modèle laisse de côté~; il n'existe aucune
forme fermée au-delà des formes asymptotiques. Il s'agit d'un problème de contrôle stochastique (livre~4, chapitre~9) dont l'objectif est une mesure de risque
du P\&L couvert.

\begin{definition}[Mesure de risque convexe, mesure de risque entropique, prix d'indifférence]\label{def:ml:deep-hedging-and-machine-learning-in-pricing:risk}
Une \emph{mesure de risque convexe}\index{mesure de risque convexe} $\rho$ associe à un P\&L aléatoire $X$ le montant de liquidités qui le rend
acceptable, elle est monotone, invariante par ajout de liquidités ($\rho(X + c) = \rho(X) - c$) et convexe (Föllmer and Schied, 2002).
L'expected shortfall (livre~6, chapitre~21) en est un. La \emph{mesure de risque entropique}\index{mesure de risque entropique} est $\rho(X) =
\lambda^{-1}\log\E[e^{-\lambda X}]$, l'équivalent certain de l'utilité exponentielle à aversion au risque $\lambda$. Le
\emph{prix d'indifférence}\index{prix d'indifférence} d'un produit est la prime qui laisse le vendeur exactement aussi bien loti, selon
la mesure de risque, que s'il ne vendait pas, chacun avec sa meilleure stratégie de trading~: $p = \inf_\delta\rho(-Z + G_\delta) -
\inf_\delta\rho(G_\delta)$, où $G_\delta$ est le gain de la stratégie de trading $\delta$ après coûts. Lorsque le
prix du sous-jacent est une martingale, comme dans la simulation du chapitre, ne pas trader est optimal sans le produit, et le second terme est nul.
\end{definition}

\section{Couverture profonde}

\begin{definition}[Couverture profonde]\label{def:ml:deep-hedging-and-machine-learning-in-pricing:dh}
La \emph{couverture profonde}\index{couverture profonde} (deep hedging) représente la stratégie de couverture par un \omterm{def:ml:neural-networks-for-noisy-tabular-data:mlp}{réseau de neurones} qui associe à
chaque date (temps, cours, positions courantes) les nouvelles positions, et l'entraîne sur des trajectoires simulées pour minimiser une \omterm{def:ml:deep-hedging-and-machine-learning-in-pricing:risk}{mesure de risque
convexe} du P\&L couvert après l'ensemble des frictions (Buehler, Gonon, Teichmann and Wood, 2019).
\end{definition}

Le réseau du chapitre (le \cref{lst:ml:dh:band}) est le réseau à bande sans transaction d'Imaki et coauteurs (2021)~: à partir du
temps jusqu'à l'échéance et la log-moneyness, il produit une bande, et la nouvelle position est l'ancienne ramenée dans celle-ci, de sorte que ne pas
trader soit le choix par défaut et que la largeur de la bande puisse être lue directement. Le P\&L de la vente du call et de la couverture quotidienne avec le
sous-jacent est calculé sur un lot de trajectoires (le \cref{lst:ml:dh:pnl}), le risque entropique avec $\lambda = 1$ sert de perte, et 1\,000
pas d'\omterm{def:ml:neural-networks-for-noisy-tabular-data:adam}{Adam} sur des lots de 2\,048 parmi 20\,000 trajectoires d'entraînement prennent environ quinze secondes sur un cœur. Chaque \omterm{def:ml:reinforcement-learning-foundations:mdp}{politique} est ensuite évaluée
sur 20\,000 autres trajectoires.

\begin{omfigure}
\begin{tikzpicture}
\begin{axis}[width=11cm,height=6cm,xlabel={coût proportionnel (points de base)},ylabel={prix d'indifférence (entropique, $\lambda = 1$)},
  tick label style={font=\small},label style={font=\small},xmin=0,xmax=20,ymin=2.4,ymax=3.05,grid=major,grid style={gray!25},
  legend cell align=left,legend style={font=\footnotesize,at={(0.5,-0.22)},anchor=north,fill=none,draw=none},
  table/col sep=comma]
\addplot[omDef,very thick,mark=*] table[x=cost,y=deep]{figdata/ml/19-deep-hedging-and-machine-learning-in-pricing/indifference.csv};
\addplot[omProp,very thick,mark=square*] table[x=cost,y=ww]{figdata/ml/19-deep-hedging-and-machine-learning-in-pricing/indifference.csv};
\addplot[omThm,very thick,mark=triangle*] table[x=cost,y=bs]{figdata/ml/19-deep-hedging-and-machine-learning-in-pricing/indifference.csv};
\legend{couverture profonde (réseau à bande),bande de Whalley--Wilmott,delta de Black--Scholes}
\end{axis}
\end{tikzpicture}

\omcaption{Le \omterm{def:ml:deep-hedging-and-machine-learning-in-pricing:risk}{prix d'indifférence} du vendeur du call à un mois à la monnaie (sa valeur de Heston est de 2.25) sous chaque
\omterm{def:ml:reinforcement-learning-foundations:mdp}{politique} de couverture, en fonction du coût de transaction du sous-jacent. Données~: \texttt{ml\_hedge.hedging}.}\label{fig:ml:dh:indiff}
\end{omfigure}

La \Cref{fig:ml:dh:indiff} illustre la première moitié du résultat phare du chapitre. Sans coûts, le \omterm{def:ml:deep-hedging-and-machine-learning-in-pricing:risk}{prix d'indifférence} de la \omterm{def:ml:deep-hedging-and-machine-learning-in-pricing:dh}{couverture profonde} est
2.454 contre 2.475 pour le delta de Black--Scholes~: sous Heston, spot et volatilité évoluent ensemble ($\rho = -0.7$) et le
réseau apprend un ratio de couverture qui l'exploite. Avec des coûts de 5, 10 et 20 points de base, il monte à 2.558, 2.658 et 2.840~; la
couverture en delta de Black--Scholes monte à 2.608, 2.743 et 3.018, et celle de la bande de Whalley--Wilmott (Whalley and Wilmott, 1997) à 2.606,
2.712 et 2.910. Le \Cref{tab:ml:dh:ten} donne les autres statistiques à 10 points de base. Entraîné plutôt sur l'expected shortfall, le
réseau abaisse l'expected shortfall (3.763 contre 3.817) et relève le risque entropique (2.671 contre 2.658)~: la
mesure de risque fait partie du modèle, et chaque réseau est optimal selon son propre critère.

\begin{table}[htbp]
\centering\small
\begin{tabular}{@{}>{\raggedright\arraybackslash}p{5.2cm}>{\centering\arraybackslash}p{2.3cm}>{\centering\arraybackslash}p{2.3cm}>{\centering\arraybackslash}p{2.3cm}>{\centering\arraybackslash}p{2.3cm}@{}}
\toprule
\omterm{def:ml:reinforcement-learning-foundations:mdp}{politique} & risque entropique & expected shortfall 95~\% & coût moyen & écart-type\\
\midrule
\omterm{def:ml:deep-hedging-and-machine-learning-in-pricing:dh}{couverture profonde}, entropique & 2.658 & 3.817 & 2.442 & 0.643\\
\omterm{def:ml:deep-hedging-and-machine-learning-in-pricing:dh}{couverture profonde}, expected shortfall & 2.671 & 3.763 & 2.457 & 0.657\\
bande de Whalley--Wilmott & 2.712 & 4.044 & 2.424 & 0.703\\
delta de Black--Scholes & 2.743 & 4.090 & 2.507 & 0.595\\
aucune couverture & 8.767 & 9.992 & 2.235 & 2.945\\
\bottomrule
\end{tabular}
\caption{Vente du call et couverture quotidienne à 10 points de base de coût, sur 20\,000 trajectoires de test~: mesures de risque du P\&L
avant la prime, sa perte moyenne et son écart-type. Données~: \texttt{ml\_hedge.hedging}, \texttt{shortfall\_hedge}.}\label{tab:ml:dh:ten}
\end{table}

La seconde moitié du résultat phare est la bande elle-même (la \cref{fig:ml:dh:bands}). À la monnaie, à mi-parcours de l'échéance, la
bande apprise a une largeur de 0.031 en delta à 5 points de base et de 0.104 à 20~; la bande de Whalley--Wilmott, $\pm(3cS\Gamma^2/2\lambda)^{1/3}$,
vaut 0.181 et 0.288. La formule asymptotique suppose un trader qui surveille en continu et négocie à la frontière de la bande~; avec une
décision par jour, la position dérive entre deux décisions et une bande large amplifie cette dérive, si bien que le réseau garde une bande plus étroite
à la monnaie et au-dessus, et une plus large hors de la monnaie, où il est presque inutile de couvrir l'option. La formule donne la bonne idée et la mauvaise taille
pour ce trader~; le réseau trouve la taille adaptée au calendrier réel du trader et à sa mesure de risque.

\begin{omfigure}
\begin{tikzpicture}
\begin{axis}[width=11cm,height=6cm,xlabel={spot (strike 100, à mi-parcours de l'échéance)},ylabel={position (actions par option)},
  tick label style={font=\small},label style={font=\small},xmin=92,xmax=108,ymin=-0.05,ymax=1.0,grid=major,grid style={gray!25},
  legend cell align=left,legend style={font=\footnotesize,at={(0.02,0.98)},anchor=north west,fill=white,draw=none},
  table/col sep=comma]
\addplot[black,thick,dashed] table[x=spot,y=delta]{figdata/ml/19-deep-hedging-and-machine-learning-in-pricing/bands.csv};
\addplot[omProp,thick] table[x=spot,y=wwlo]{figdata/ml/19-deep-hedging-and-machine-learning-in-pricing/bands.csv};
\addplot[omDef,very thick] table[x=spot,y=deeplo]{figdata/ml/19-deep-hedging-and-machine-learning-in-pricing/bands.csv};
\addplot[omProp,thick,forget plot] table[x=spot,y=wwhi]{figdata/ml/19-deep-hedging-and-machine-learning-in-pricing/bands.csv};
\addplot[omDef,very thick,forget plot] table[x=spot,y=deephi]{figdata/ml/19-deep-hedging-and-machine-learning-in-pricing/bands.csv};
\legend{delta de Black--Scholes,bande de Whalley--Wilmott,bande apprise}
\end{axis}
\end{tikzpicture}

\omcaption{Bandes sans transaction à 20 points de base de coût, à mi-parcours de l'échéance~: la position est laissée inchangée à l'intérieur d'une bande et ramenée
au bord le plus proche si elle se trouve à l'extérieur. Données~: \texttt{ml\_hedge} (\texttt{fig\_hedge.py}).}\label{fig:ml:dh:bands}
\end{omfigure}

\section{Pricers de substitution et apprentissage différentiel}

\begin{definition}[Modèle de substitution, apprentissage automatique différentiel]\label{def:ml:deep-hedging-and-machine-learning-in-pricing:surrogate}
Un \emph{modèle de substitution}\index{modèle de substitution} (surrogate model) est une approximation rapide d'une fonction lente, ici un pricer, ajustée à ses
sorties sur des entrées échantillonnées. L'\emph{apprentissage automatique différentiel}\index{apprentissage automatique différentiel} (differential machine learning) entraîne le modèle de substitution sur les
dérivées des étiquettes par rapport aux entrées, en plus des étiquettes elles-mêmes, les dérivées propres du réseau étant calculées par
différentiation automatique (livre~4, chapitre~28)~; avec des étiquettes de Monte Carlo, les étiquettes de dérivées sont pathwise (Huge et Savine,
2020).
\end{definition}

Un système de gestion des risques qui valorise un portefeuille sous des milliers de scénarios ne peut pas se permettre une intégrale de Fourier ni une simulation de Monte Carlo pour chaque
prix. Le \omterm{def:ml:deep-hedging-and-machine-learning-in-pricing:surrogate}{modèle de substitution} du chapitre apprend le prix de Heston d'une option d'achat à trois mois en fonction du spot (70 à 130) et de la
variance initiale (0.01 à 0.09), à partir des étiquettes les moins coûteuses qui soient~: une trajectoire simulée par échantillon, dont le payoff est un prix sans biais
mais très bruité. L'étiquette différentielle est la dérivée pathwise du payoff par rapport au spot,
$\mathbf1\{S_T > K\}S_T/S_0$, gratuite une fois la trajectoire obtenue. Le \Cref{tab:ml:dh:surr} confronte les deux apprentissages au prix de Fourier
sur 400 points de test (\cref{lst:ml:dh:surr} est la boucle d'apprentissage).

\begin{table}[htbp]
\centering\small
\begin{tabular}{@{}rcccc@{}}
\toprule
& \multicolumn{2}{c}{RMSE du prix} & \multicolumn{2}{c}{RMSE du delta}\\
\cmidrule(lr){2-3}\cmidrule(l){4-5}
échantillons d'apprentissage & standard & différentiel & standard & différentiel\\
\midrule
256 & 0.927 & 0.679 & 0.066 & 0.057\\
1\,024 & 0.436 & 0.218 & 0.059 & 0.021\\
4\,096 & 0.297 & 0.279 & 0.042 & 0.020\\
\bottomrule
\end{tabular}
\caption{\omterm{def:ml:deep-hedging-and-machine-learning-in-pricing:surrogate}{Modèle de substitution} à deux couches du prix de Heston d'une option d'achat, appris à partir de payoffs de Monte Carlo à une seule trajectoire, avec et sans
étiquettes de delta pathwise~: erreurs par rapport au prix de Fourier et à son delta par différences finies sur 400 points de test. Données~:
\texttt{ml\_hedge.surrogates}.}\label{tab:ml:dh:surr}
\end{table}

Avec 1\,024 échantillons, les étiquettes de dérivées divisent par deux l'erreur sur le prix et réduisent de près des deux tiers l'erreur sur le delta~; avec 4\,096, les
prix sont aussi bons dans les deux cas, mais seuls les deltas du réseau différentiel le sont. Le delta d'un \omterm{def:ml:deep-hedging-and-machine-learning-in-pricing:surrogate}{modèle de substitution} est celui qu'utilise le hedger, et
un réseau ajusté aux seuls prix n'a aucune raison de bien estimer sa pente.

\section{Réseaux de calibration}

\begin{definition}[Calibration profonde]\label{def:ml:deep-hedging-and-machine-learning-in-pricing:calib}
La \emph{calibration profonde}\index{calibration profonde} (deep calibration) remplace l'optimisation numérique des paramètres d'un modèle face aux cotations de marché
(livre~4, chapitre~24) par un réseau entraîné sur des données générées par le modèle~: soit l'application inverse des cotations vers les paramètres (Hernandez,
2016), soit une application directe rapide des paramètres vers les cotations, insérée dans un optimiseur standard (Horvath, Muguruza et Tomas, 2021).
\end{definition}

Le réseau du chapitre apprend l'application inverse, qui va de quinze prix d'options d'achat normalisés (cinq strikes, trois échéances) vers les cinq
paramètres de Heston, sur 4\,000 jeux de paramètres tirés uniformément dans un pavé. Sur 200 jeux de test, ses erreurs, rapportées à l'amplitude de chaque
paramètre, sont de 1.7~\% pour la variance initiale, 18.7~\% pour la vitesse de retour à la moyenne, 9.7~\% pour la variance de long terme,
9.0~\% pour la volatilité de la variance et 9.8~\% pour la corrélation~; les surfaces revalorisées avec ses réponses s'écartent des vraies
de 13.2 points de base de spot en moyenne. La vitesse et la variance de long terme sont mal identifiées par trois échéances, et le réseau
est aussi incertain à leur sujet que ne le sont les données. Sur cinq surfaces de test, sa réponse s'écarte des volatilités implicites de 0.46 point en
moyenne, et le calibrateur de Levenberg--Marquardt du livre~5, initialisé à partir de cette réponse ou d'une estimation fixe, les ajuste exactement. Le réseau est
une première estimation rapide et un contrôle de cohérence, non un remplacement, sur des données propres issues du modèle d'entraînement~; sur des données de marché,
que nul modèle n'ajuste exactement, il n'a jamais vu les résidus qu'il rencontrera.

\section{Où le pricing appris échoue}

Trois défaillances reviennent. \emph{Extrapolation}~: un \omterm{def:ml:deep-hedging-and-machine-learning-in-pricing:surrogate}{modèle de substitution} ou un réseau de calibration situé hors de son pavé d'entraînement renvoie des chiffres
assurés, sans aucune alerte, précisément là où un marché sous tension emporte le portefeuille~; entraîner sur des pavés plus larges que tout scénario et vérifier
les entrées par rapport à eux. \emph{Arbitrage}~: les prix d'un \omterm{def:ml:deep-hedging-and-machine-learning-in-pricing:surrogate}{modèle de substitution} ne sont pas nécessairement monotones ou convexes en fonction du strike, et ses grecques peuvent
changer de signe là où elles ne le devraient pas~; contraindre l'architecture ou vérifier les sorties. \emph{Risque de modèle}~: une \omterm{def:ml:deep-hedging-and-machine-learning-in-pricing:dh}{couverture profonde} est optimale
pour le simulateur sur lequel elle a été entraînée, comme l'agent du chapitre~17~; l'entraîner sur plusieurs modèles ou sur des paramètres aléatoires, et
la comparer à la bande asymptotique sur des trajectoires qui n'ont pas servi à l'entraîner.

\begin{method}[Apprentissage automatique sur un desk de produits dérivés]\label{met:ml:dh:method}
\begin{enumerate}
\item Conserver le pricer exact comme référence~; chaque \omterm{def:ml:deep-hedging-and-machine-learning-in-pricing:surrogate}{modèle de substitution} est validé par rapport à lui sur une grille plus large que les
scénarios du portefeuille, pour les prix comme pour les grecques.
\item Entraîner des \omterm{def:ml:deep-hedging-and-machine-learning-in-pricing:surrogate}{modèles de substitution} avec des étiquettes différentielles lorsque les dérivées sont peu coûteuses (pathwise ou par différentiation automatique).
\item Utiliser les réseaux de calibration comme points de départ et comme contrôles~; l'optimiseur a le dernier mot.
\item Entraîner des \omterm{def:ml:deep-hedging-and-machine-learning-in-pricing:dh}{couvertures profondes} sur les frictions et la mesure de risque propres au desk, et comparer à la bande en forme fermée hors échantillon et
à travers les modèles.
\end{enumerate}
\end{method}

\section{Tutoriel~: couverture en présence de frictions}\label{tut:ml:deep-hedging-and-machine-learning-in-pricing:tut}

\begin{tutorial}
\textbf{Objectif.} Appliquer la \omterm{def:ml:deep-hedging-and-machine-learning-in-pricing:dh}{couverture profonde} à l'option d'achat à quatre niveaux de coût et relever les \omterm{def:ml:deep-hedging-and-machine-learning-in-pricing:risk}{prix d'indifférence} et les bandes~; entraîner le \omterm{def:ml:deep-hedging-and-machine-learning-in-pricing:surrogate}{modèle de substitution} avec
et sans étiquettes différentielles~; ajuster le réseau de calibration. \textbf{État final~:}
les \cref{fig:ml:dh:indiff,fig:ml:dh:bands}, \cref{tab:ml:dh:ten,tab:ml:dh:surr}.
\begin{tutsteps}
\item \textbf{Le réseau de bandes sans transaction.}
\omcode{code/firm/deephedge/firm_deephedge.py}{54}{69}{Un réseau qui produit une bande~; la position y est ramenée.}\label{lst:ml:dh:band}
\item \textbf{Le P\&L de couverture après coûts.}
\omcode{code/firm/deephedge/firm_deephedge.py}{72}{86}{Vente d'une option d'achat et sa couverture sur un lot de trajectoires.}\label{lst:ml:dh:pnl}
\item \textbf{Apprentissage différentiel.}
\omcode{code/firm/deephedge/firm_deephedge.py}{165}{186}{Un modèle de substitution entraîné sur les prix et les dérivées pathwise.}\label{lst:ml:dh:surr}
\item \textbf{Exécuter} \texttt{ml\_hedge.hedging()} (environ une minute et demie sur un cœur), \texttt{surrogates()},
\texttt{calibration()}, \texttt{lm\_check()} et \texttt{fig\_hedge.py}.
\end{tutsteps}
\textbf{Pour aller plus loin.} Donner au réseau la variance courante comme entrée et ajouter un swap de variance comme seconde couverture~;
entraîner sur des paramètres de Heston aléatoires et tester sur un autre modèle.
\end{tutorial}

\section{Construction~: couverture profonde et pricers appris}\label{bld:ml:deep-hedging-and-machine-learning-in-pricing:deephedge}

\begin{build}
\bfield{Rôle} Couvertures et prix appris là où ils font gagner du temps ou des coûts, vérifiés par rapport aux moteurs exacts du livre~5.
\bfield{Interface} \texttt{heston\_paths}~; \texttt{BandNet}, \texttt{hedge\_pnl(policy, S, K, T, cost)}, \texttt{entropic},
\texttt{expected\_shortfall}, \texttt{train\_hedge(S, K, T, cost, risk, seed, steps, batch, lr)}~; \texttt{bs\_policy},
\texttt{ww\_policy}, \texttt{no\_hedge}, \texttt{band}~; \texttt{SurrogateNet},
\texttt{fit\_surrogate(X, y, dy, differential, seed)}, \texttt{surrogate\_values}~; \texttt{fit\_calibrator(vols, params, seed)}.
\bfield{Règles} Chaque prix, grecque et couverture appris est évalué par rapport à une référence exacte sur des données qui n'ont pas servi à l'entraînement~;
les pavés d'entraînement et les niveaux de coût sont consignés avec le modèle.
\bfield{Tests d'acceptation} \texttt{code/firm/deephedge/tests/}~: le P\&L de la couverture en delta sans coûts est du même ordre que l'erreur
de couverture discrète~; la \omterm{def:ml:deep-hedging-and-machine-learning-in-pricing:risk}{mesure de risque entropique} d'une constante est l'opposé de cette constante~; le réseau de bandes maintient sa position à l'intérieur de la
bande~; la bande de Whalley--Wilmott s'élargit avec le coût~; un \omterm{def:ml:deep-hedging-and-machine-learning-in-pricing:surrogate}{modèle de substitution} différentiel d'une fonction connue retrouve sa pente~; le
calibrateur inverse une application simple.
\bfield{Pour aller plus loin} Un second instrument de couverture~; \omterm{def:ml:reinforcement-learning-foundations:mdp}{politiques} récurrentes avec la variance comme état caché~; architectures de
\omterm{def:ml:deep-hedging-and-machine-learning-in-pricing:surrogate}{modèles de substitution} sans arbitrage.
\end{build}

\begin{omsources}
\item H.~Buehler, L.~Gonon, J.~Teichmann et B.~Wood, ``Deep hedging'', \emph{Quantitative Finance} 19(8), 2019.
\item S.~Imaki, K.~Imajo, K.~Ito, K.~Minami et K.~Nakagawa, ``No-transaction band network'', arXiv:2103.01775, 2021.
\item A.~E.~Whalley et P.~Wilmott, ``An asymptotic analysis of an optimal hedging model for option pricing with transaction
costs'', \emph{Mathematical Finance} 7(3), 1997.
\item H.~Föllmer et A.~Schied, ``Convex measures of risk and trading constraints'', \emph{Finance and Stochastics} 6(4), 2002.
\item B.~Huge et A.~Savine, ``Differential machine learning'', arXiv:2005.02347, 2020.
\item B.~Horvath, A.~Muguruza et M.~Tomas, ``Deep learning volatility'', \emph{Quantitative Finance} 21(1), 2021.
\item A.~Hernandez, ``Model calibration with neural networks'', SSRN 2812140, 2016.
\end{omsources}

\section{Exercices}

\begin{exercise}[$\star$]\label{exo:ml:deep-hedging-and-machine-learning-in-pricing:1}
Montrer que la \omterm{def:ml:deep-hedging-and-machine-learning-in-pricing:risk}{mesure de risque entropique} est invariante par cash~: $\rho(X + c) = \rho(X) - c$. Que vaut $\rho$ d'une constante~?
\end{exercise}

\begin{exercise}[$\star$]\label{exo:ml:deep-hedging-and-machine-learning-in-pricing:2}
Calculer la demi-largeur de Whalley--Wilmott à la monnaie, à mi-chemin de l'échéance, pour $c = 0.001$, $S = 100$, $\lambda = 1$ et un
gamma de Black--Scholes égal à 0.0998.
\end{exercise}

\begin{exercise}[$\star$]\label{exo:ml:deep-hedging-and-machine-learning-in-pricing:3}
Pourquoi le \omterm{def:ml:deep-hedging-and-machine-learning-in-pricing:risk}{prix d'indifférence} est-il supérieur à la valeur de Heston (2.25) même en l'absence de coûts~?
\end{exercise}

\begin{exercise}[$\star\star$]\label{exo:ml:deep-hedging-and-machine-learning-in-pricing:4}
Pourquoi le deep hedging bat-il le delta de Black--Scholes lorsqu'il n'y a aucun coût du tout~?
\end{exercise}

\begin{exercise}[$\star\star$]\label{exo:ml:deep-hedging-and-machine-learning-in-pricing:5}
Pourquoi la vitesse de retour à la moyenne et la variance de long terme sont-elles difficiles à apprendre pour le réseau de calibration, et quelles données aideraient~?
\end{exercise}

\begin{exercise}[$\star\star$]\label{exo:ml:deep-hedging-and-machine-learning-in-pricing:6}
\emph{Trouver la faille.} «~Notre \omterm{def:ml:deep-hedging-and-machine-learning-in-pricing:surrogate}{modèle de substitution} reproduit le pricer à 0.1~\% près sur l'\omterm{def:ml:why-financial-machine-learning-is-different:sets}{ensemble de test}, nous l'utilisons donc pour les scénarios de stress.~»
\end{exercise}

\begin{exercise}[$\star\star\star$]\label{exo:ml:deep-hedging-and-machine-learning-in-pricing:7}
\emph{Programmation.} Dériver l'étiquette de delta par la méthode des trajectoires pour un call digital ($\mathbf1\{S_T > K\}$). Qu'est-ce qui
ne fonctionne pas, et que fait l'approche de Huge et Savine à la place~?
\end{exercise}

\begin{exercise}[$\star\star\star$]\label{exo:ml:deep-hedging-and-machine-learning-in-pricing:8}
Montrer que, sans coûts et dans un marché complet (Black--Scholes, négociation en continu), le \omterm{def:ml:deep-hedging-and-machine-learning-in-pricing:risk}{prix
d'indifférence} entropique est égal au prix de Black--Scholes pour tout $\lambda$.
\end{exercise}

\section{Problème~: couverture avec frictions}

\begin{problem}[{Problème du week-end --- quand ne pas trader}]\label{pb:ml:deep-hedging-and-machine-learning-in-pricing:1}
Le call du chapitre, les couvertures, le \omterm{def:ml:deep-hedging-and-machine-learning-in-pricing:surrogate}{modèle de substitution} et le calibrateur.

\medskip\noindent\textbf{Partie I --- Le problème.}
\begin{enumerate}
\item Qu'est-ce qui est vendu, comment est-ce couvert, et que voit le réseau~?
\item Combien coûte la couverture en delta quotidienne à 5 points de base, et comment cela se compare-t-il au vega~?
\item Pourquoi un réseau de bandes plutôt qu'un réseau qui produit directement la position~?
\item Quelle est la perte, et que signifie $\lambda = 1$ ici~?
\end{enumerate}

\medskip\noindent\textbf{Partie II --- Résultats.}
\begin{enumerate}[resume]
\item Donner les \omterm{def:ml:deep-hedging-and-machine-learning-in-pricing:risk}{prix d'indifférence} des trois couvertures à chaque niveau de coût.
\item Comparer la bande apprise à la bande de Whalley--Wilmott.
\item Qu'est-ce que l'entraînement sur l'expected shortfall change~?
\item Pourquoi la couverture de Whalley--Wilmott a-t-elle le coût moyen le plus faible à 10 points de base, mais pas le risque le plus faible~?
\end{enumerate}

\medskip\noindent\textbf{Partie III --- \omterm{def:ml:deep-hedging-and-machine-learning-in-pricing:surrogate}{Modèles de substitution} et calibration.}
\begin{enumerate}[resume]
\item À quoi servent les étiquettes différentielles, et quand~?
\item Quelles erreurs commet le réseau de calibration~?
\item Comment faut-il combiner le réseau et l'optimiseur~?
\item Où chaque pricer appris échouerait-il sur un desk réel~?
\end{enumerate}

\medskip\noindent\textbf{Partie IV --- Le verdict.}
\begin{enumerate}[resume]
\item Énoncer le \emph{résultat nommé}~: le \omterm{def:ml:deep-hedging-and-machine-learning-in-pricing:risk}{prix d'indifférence} en fonction du niveau de coût, et la largeur de la bande de non-négociation apprise
comparée à celle de la bande de Whalley--Wilmott.
\item Faut-il coter l'option au \omterm{def:ml:deep-hedging-and-machine-learning-in-pricing:risk}{prix d'indifférence}~? Pourquoi~?
\item Comment valider une \omterm{def:ml:deep-hedging-and-machine-learning-in-pricing:dh}{couverture profonde} avant de l'utiliser sur un book~?
\item Que changerait un second instrument de couverture~?
\item Comment garder un \omterm{def:ml:deep-hedging-and-machine-learning-in-pricing:surrogate}{modèle de substitution} fiable lorsque le modèle est recalibré chaque jour~?
\item Quel est le risque d'entraîner la couverture sur un seul jeu de paramètres de Heston~?
\item Où la différentiation automatique intervient-elle dans ce chapitre~?
\item En une phrase~: qu'apprend un réseau de couverture que le delta n'apprend pas~?
\end{enumerate}
\end{problem}

\section{Questions d'entretien}

\begin{interviewq}[$\star$ \iqroles{chercheur, trader}]\label{iq:ml:deep-hedging-and-machine-learning-in-pricing:1}
Pourquoi la couverture en delta avec coûts de transaction appelle-t-elle une bande de non-négociation~?
\end{interviewq}

\begin{interviewq}[$\star\star$ \iqroles{chercheur}]\label{iq:ml:deep-hedging-and-machine-learning-in-pricing:2}
Qu'est-ce que le deep hedging, et quelles sont ses entrées, ses sorties et sa perte~?
\end{interviewq}

\begin{interviewq}[$\star\star$ \iqroles{ingénieur ML, chercheur}]\label{iq:ml:deep-hedging-and-machine-learning-in-pricing:3}
Expliquez l'\omterm{def:ml:deep-hedging-and-machine-learning-in-pricing:surrogate}{apprentissage automatique différentiel} et pourquoi il aide avec des étiquettes Monte Carlo bruitées.
\end{interviewq}

\begin{interviewq}[$\star\star$ \iqroles{chercheur}]\label{iq:ml:deep-hedging-and-machine-learning-in-pricing:4}
Qu'est-ce qu'un \omterm{def:ml:deep-hedging-and-machine-learning-in-pricing:risk}{prix d'indifférence}, et quel est son lien avec une mesure de risque~?
\end{interviewq}

\begin{interviewq}[$\star\star$ \iqroles{ingénieur ML}]\label{iq:ml:deep-hedging-and-machine-learning-in-pricing:5}
Un réseau de calibration est rapide. Pourquoi ne pas l'utiliser seul~?
\end{interviewq}

\begin{interviewq}[$\star\star\star$ \iqroles{chercheur}]\label{iq:ml:deep-hedging-and-machine-learning-in-pricing:6}
Le gamma de votre \omterm{def:ml:deep-hedging-and-machine-learning-in-pricing:surrogate}{modèle de substitution} devient négatif pour des strikes profondément hors de la monnaie. Qu'est-il arrivé et comment le corrigez-vous~?
\end{interviewq}
